\chapter*{Annexes}
\addcontentsline{toc}{chapter}{Annexes}
\markboth{Annexes}{}
\stepcounter{chapter}
\addtocontents{lot}{\vspace{3.8mm}}
\addtocontents{lof}{\vspace{3.8mm}}

Ces annexes rassemblent les éléments trop détaillés pour figurer dans le corps du mémoire, mais
nécessaires à qui veut vérifier une affirmation ou reprendre le travail. Toutes les valeurs qui y
figurent proviennent du dépôt du projet, et aucune n'a été recalculée pour les besoins de la
rédaction.

%===================== ANNEXE A =====================%
\section*{Annexe A : registre des risques}
\addcontentsline{toc}{section}{Annexe A : registre des risques}
\renewcommand{\thetable}{A.\arabic{table}}
\setcounter{table}{0}

Ce registre recense les risques rencontrés pendant le projet, ce qu'ils menaçaient et l'état dans
lequel le mémoire les laisse. Il est tenu à jour dans le dépôt. Un risque est clos lorsque le
défaut a été corrigé et la correction vérifiée ; il reste ouvert lorsque la limite est connue et
assumée.

\begin{center}
\styletableau
\begin{longtable}{|Cp{1.2cm}|Lp{10.6cm}|Cp{2.0cm}|}
\caption{Registre des risques du projet}\label{tab:registre}\\
\hline
\rowcolor{enteteTab}
\textbf{ID} & \textbf{Risque} & \textbf{Statut} \\
\hline
\endfirsthead
\hline
\rowcolor{enteteTab}
\textbf{ID} & \textbf{Risque} & \textbf{Statut} \\
\hline
\endhead
R-01 & Redondance du corpus : peu de types de tickets distincts, d'où un risque de fuite et une
généralisation limitée & Ouvert \\
\hline
R-02 & Sur-apprentissage : écart entre la validation croisée et le jeu de test & Ouvert \\
\hline
R-03 & Rappel faible sur la catégorie \emph{Sécurité}, en deçà de la cible & Ouvert \\
\hline
R-04 & Hallucinations du modèle de langage dans la synthèse : incident reproduit et corrigé par
deux garde-fous, la vérification de fidélité restant non automatisée & Ouvert \\
\hline
R-05 & Conformité au règlement sur les données personnelles : anonymisation et rétention à
formaliser & Ouvert \\
\hline
R-06 & Périmètre trop large pour la durée du stage & Ouvert \\
\hline
R-07 & Dette d'intégration entre l'interface et le service & Clos \\
\hline
R-08 & Chemin d'ingestion externe non fonctionnel & Clos \\
\hline
R-09 & Secrets de déploiement laissés à leur valeur par défaut dans les fichiers d'exemple &
Clos \\
\hline
R-10 & Coût et durée d'accès au matériel pour l'adaptation du transformer & Ouvert \\
\hline
R-11 & Dérive documentaire entre le code et les documents de pilotage & Ouvert \\
\hline
R-12 & Panne de démonstration le jour de la soutenance & Ouvert \\
\hline
R-13 & Temps de réponse hors budget, et dégradation forte sous concurrence & Ouvert \\
\hline
R-14 & Jeu de test non représentatif : trois catégories sur huit n'y sont représentées par aucun
exemple & Ouvert \\
\hline
R-15 & Avis du demandeur affiché comme enregistré sans être transmis au service & Clos \\
\hline
R-16 & Aucun routage différencié selon la confiance : l'indice est calculé sans piloter aucune
automatisation & Ouvert \\
\hline
R-17 & Courses entre processus lors de la migration vers la base relationnelle et la base
vectorielle & Clos \\
\hline
R-18 & Couverture de test absente sur les fonctionnalités développées entre juillet et août &
Ouvert \\
\hline
R-19 & Latence de l'inférence du modèle de langage sur processeur & Ouvert \\
\hline
R-20 & Seuils calibrés sur le jeu d'évaluation et non sur un volume observé, faute de trafic
réel : déclencheurs de réentraînement, et arbitrage du seuil de pertinence entre pureté et nombre
des suggestions & Ouvert \\
\hline
R-21 & Fuite par paraphrase non détectée à la construction du jeu d'apprentissage & Ouvert \\
\hline
R-22 & Annotation de référence portée par l'auteur et non par un consultant du service : elle établit qu'une résolution traite le problème, non qu'elle aurait fait gagner du temps &
Ouvert \\
\hline
R-23 & Règles métier ajustées sur le jeu de test : quatre mots-clés retirés après audit & Clos \\
\hline
R-24 & Aucun référentiel de mots-clés validé formellement par le service support & Ouvert \\
\hline
R-25 & Coexistence de mesures anciennes et récentes du moteur de recherche, sur des index
différents & Ouvert \\
\hline
R-26 & Informations d'entreprise manquantes, dont l'autorisation d'exploiter les données réelles &
Ouvert \\
\hline
\end{longtable}
\end{center}

\FloatBarrier

%===================== ANNEXE B =====================%
\section*{Annexe B : cycle de vie du ticket}
\addcontentsline{toc}{section}{Annexe B : cycle de vie du ticket}
\renewcommand{\thetable}{B.\arabic{table}}
\setcounter{table}{0}

Le diagramme d'états du chapitre~7 donne la vue d'ensemble ; cette annexe en donne le détail. Le
tableau qui suit énumère les transitions que le service autorise réellement, telles que la table de
contrôle du code les déclare. Toute transition qui n'y figure pas est refusée par le serveur, quoi
que propose l'interface.

\begin{table}[!htbp]
\centering
\styletableau
\caption{Transitions d'état autorisées}
\label{tab:transitions}
\begin{tabular}{|Bp{3.6cm}|Lp{9.4cm}|}
\hline
\rowcolor{enteteTab}
\textbf{État courant} & \textbf{États atteignables} \\
\hline
En attente & Prise en charge, Rejeté \\
\hline
Prise en charge & Résolu, Problème, Rejeté, En attente \\
\hline
Problème & Prise en charge, Résolu, Rejeté \\
\hline
Résolu & Aucun \\
\hline
Rejeté & Aucun \\
\hline
\end{tabular}
\end{table}

Les états \emph{Résolu} et \emph{Rejeté} n'ont aucune transition sortante, sans être figés pour
autant. La catégorie, la priorité et la solution appliquée restent modifiables après clôture, le
contrôle n'exigeant une transition permise que lorsque le statut soumis diffère du statut courant.

\FloatBarrier

%===================== ANNEXE C =====================%
% L'ancienne annexe D reprenait la distribution complete des latences. Elle
% doublait la figure de la campagne de charge du chapitre 6, qui porte les
% memes releves ; retiree le 2026-10-14, les annexes suivantes avancent d'une
% lettre.
\section*{Annexe C : paramètres de la chaîne de traitement}
\addcontentsline{toc}{section}{Annexe C : paramètres de la chaîne de traitement}
\renewcommand{\thetable}{C.\arabic{table}}
\setcounter{table}{0}

Les valeurs qui suivent gouvernent le comportement du système. Elles tiennent dans un fichier
unique, versionné avec le code : une exécution de la chaîne est donc entièrement décrite par son
contenu, et deux exécutions d'un même fichier donnent le même résultat.

\begin{table}[!htbp]
\centering
\styletableau
\caption{Paramètres gouvernant la chaîne de traitement}
\label{tab:parametres}
\begin{tabular}{|Bp{3.2cm}|Lp{5.2cm}|Cp{2.4cm}|}
\hline
\rowcolor{enteteTab}
\textbf{Domaine} & \textbf{Paramètre} & \textbf{Valeur} \\
\hline
Découpage & Part réservée au test & 0,20 \\
\hline
Découpage & Regroupement par titre & Activé \\
\hline
Découpage & Jeu de test gelé & Activé \\
\hline
Représentation & Nombre maximal de termes & 5\,000 \\
\hline
Représentation & Longueur des séquences de mots & 1 à 2 \\
\hline
Recherche & Nombre de passages retournés & 3 \\
\hline
Recherche & Seuil de pertinence & 0,35 \\
\hline
Recherche & Seuil de repli & 0,20 \\
\hline
Recherche & Pondération de la part sémantique & 0,50 \\
\hline
Classification & Seuil de confiance & 0,60 \\
\hline
Promotion & Gain minimum exigé & 0,005 \\
\hline
Promotion & Régression maximale par catégorie & 0,05 \\
\hline
Promotion & Régression de l'exactitude & Interdite \\
\hline
Réentraînement & Tickets nouveaux requis & 150 \\
\hline
Réentraînement & Tickets requis par catégorie minoritaire & 30 \\
\hline
Réentraînement & Délai minimal entre deux promotions & 7 jours \\
\hline
Collecte & Longueur minimale d'une description & 5 mots \\
\hline
Collecte & Étiquettes faibles retenues & Non \\
\hline
Collecte & Poids d'une étiquette corrigée & 1,0 \\
\hline
Collecte & Poids d'une étiquette confirmée & 0,7 \\
\hline
\end{tabular}
\end{table}

Deux de ces valeurs appellent une réserve. Le seuil de pertinence valait 0,50, calibré sur l'échelle
du moteur sémantique et conservé lorsque le moteur hybride l'a remplacé, alors que la fusion change
l'échelle des scores : il était devenu inatteignable. Il a été ramené à 0,35 sur la distribution
mesurée, décrite à la section~\ref{sec:seuil5} ; cette calibration repose sur les quatre-vingts
requêtes du protocole d'évaluation, non sur du trafic réel. Les seuils de réentraînement, eux,
restent des ordres de grandeur raisonnés, faute d'un volume observé sur lequel les calibrer.

\FloatBarrier

%===================== ANNEXE D =====================%
\section*{Annexe D : traçabilité des chiffres publiés}
\addcontentsline{toc}{section}{Annexe D : traçabilité des chiffres publiés}
\renewcommand{\thetable}{D.\arabic{table}}
\setcounter{table}{0}

Un même indicateur peut prendre deux valeurs dans ce mémoire sans qu'aucune soit fausse. Le
F1-macro de la chaîne lexicale vaut 0,5827 d'un côté et 0,4372 de l'autre : ce ne sont ni le même
jeu de test, ni le même découpage, ni le même nombre de catégories. Les rapprocher serait une
faute, et c'est pour l'éviter que cette annexe existe. Elle donne, pour chaque valeur publiée, la
configuration mesurée, le protocole employé et le fichier du dépôt qui la porte. Les
tableaux~\ref{tab:croise}, \ref{tab:prod}, \ref{tab:moteurs5} et~\ref{tab:baseline} y renvoient.

Le \emph{protocole des expériences} sert à comparer les approches entre elles. Son découpage est
tiré par catégorie sur les 58 titres, et son jeu de test de 355 tickets représente les
8 catégories. Le \emph{protocole du pipeline} mesure le modèle réellement servi. Son découpage est
tiré par titres, son jeu de test est gelé, et il ne compte que 356 tickets pour 5 catégories. Le
tableau~\ref{tab:protocoles} les met face à face.

La même distinction vaut pour les modèles. Le \emph{meilleur modèle mesuré} est un CamemBERT adapté
par LoRA, complété des règles métier ; il n'est pas servi. Le \emph{modèle servi} est la chaîne
lexicale TF-IDF, complétée des mêmes règles. Le tableau~\ref{tab:cout} en donne le prix.

Les fichiers cités se trouvent dans \texttt{helpdesk-ai/metrics/} et l'extension \texttt{.json} est
sous-entendue. Deux exceptions : \texttt{couverture\_tests.txt}, et \texttt{feuille\_jugement}, qui
est dans \texttt{helpdesk-ai/scripts/verification/}. Cinq valeurs n'ont pas de fichier parce
qu'elles ne sont pas mesurées : les quatre écarts du tableau~\ref{tab:croise} se calculent entre
deux de ses lignes, et le 0,5719 de la section~\ref{sec:pourquoi} est un objectif à atteindre, pas
un relevé.

\begin{longtable}{|Cp{1.7cm}|Bp{3.9cm}|Lp{3.1cm}|Lp{4.0cm}|}
\caption{Origine de chaque chiffre publié}
\label{tab:chiffres} \\
\hline
\rowcolor{enteteTab}
\textbf{Valeur} & \textbf{Ce qu'elle mesure} & \textbf{Protocole} & \textbf{Fichier du dépôt} \\
\hline
\endfirsthead
\hline
\rowcolor{enteteTab}
\textbf{Valeur} & \textbf{Ce qu'elle mesure} & \textbf{Protocole} & \textbf{Fichier du dépôt} \\
\hline
\endhead

0,1382 & F1-macro, tirage au hasard stratifié & Expériences, 355 tickets, 8 catégories &
\texttt{campagne\_\allowbreak apres\_\allowbreak nettoyage} \\
\hline
0,3753 & F1-macro, TF-IDF seul & Expériences & idem \\
\hline
0,5827 & F1-macro, TF-IDF + règles & Expériences & idem \\
\hline
0,5685 & F1-macro, CamemBERT + LoRA seul & Expériences & idem \\
\hline
\textbf{0,8343} & \textbf{F1-macro, CamemBERT + LoRA + règles — meilleur modèle mesuré, non servi} &
Expériences & idem \\
\hline
\textbf{0,8644} / 0,90 & \textbf{Le même, mesuré sur l'entrée de production} : F1-macro et rappel
\emph{Sécurité} & Expériences, 355 tickets &
\texttt{camembert\_\allowbreak servable} \\
\hline
+0,0303 & Écart médian entre la configuration publiée et l'entrée de production, sur dix
découpages appariés & Expériences &
\texttt{camembert\_\allowbreak servable\_\allowbreak robustesse} \\
\hline
5,0 / 325,8 / 58,7\,ms & Coût d'inférence par ticket : TF-IDF, CamemBERT sur processeur, sur carte
graphique & Expériences & idem, clé \texttt{cout\_de\_service} \\
\hline
128,5 / 305,9\,ms & Coût de CamemBERT par ticket à 48 puis 128 jetons, qualité inchangée &
Expériences & \texttt{cout\_\allowbreak longueur\_\allowbreak sequence} \\
\hline
0,2031 & F1-macro, modèle servi sans les règles & Pipeline, 356 tickets, 8 catégories dont 3 sans
ticket & \texttt{ml\_eval} \\
\hline
\textbf{0,4372} & \textbf{F1-macro, modèle servi — chaîne lexicale + règles, référentiel entier} &
Pipeline & idem \\
\hline
0,6995 & F1-macro sur les 5 catégories représentées, métrique de décision de la promotion &
Pipeline & idem \\
\hline
0,8202 & Exactitude du modèle servi & Pipeline & idem \\
\hline
37,4\,\% & Part des tickets requalifiés par une règle & Pipeline & idem \\
\hline
0,2112 → 0,2031 & F1-macro du modèle seul, description seule puis titre + description &
Pipeline & \texttt{coherence\_\allowbreak entree\_\allowbreak modele} \\
\hline
0,4294 → 0,4372 & F1-macro avec règles, mêmes deux variantes & Pipeline & idem \\
\hline
94\,\% → 80\,\% & Part des tickets sous le seuil de confiance de 0,60 & Pipeline & idem \\
\hline
0,2417 / 0,3333 / 0,3667 & Précision@3 sans filtre : lexical, sémantique, hybride &
Évaluation RAG, 40 requêtes annotées, jugement humain & \texttt{rag\_\allowbreak seuil\_\allowbreak balayage} \\
\hline
0,5185 & Précision de l'affiché, moteur hybride au seuil servi & Évaluation RAG &
\texttt{rag\_eval} \\
\hline
0,3366 à 0,4886 & Plage du meilleur score hybride par requête, qui a commandé le seuil &
Évaluation RAG, 40 requêtes annotées &
\texttt{rag\_\allowbreak seuil\_\allowbreak balayage} \\
\hline
0,4463 / 0,5719 & F1 de \emph{Microsoft} mesuré, et valeur qu'il faudrait pour atteindre ENF-02 &
Pipeline, valeur calculée & \texttt{ml\_eval} \\
\hline
330\,ms & Temps de réponse au 95\textsuperscript{e} centile & Campagne de charge &
\texttt{load\_test} \\
\hline
260 / 234 & Tests automatisés, et sous-ensemble rejoué par l'intégration continue & — &
\texttt{suite\_\allowbreak tests.json} \\
\hline
51\,\% / 58\,\% & Couverture : code de traitement seul, puis service entier & — &
\texttt{couverture\_\allowbreak tests.txt} \\
\hline
0,3500 / 0,5333 / 0,5678 & Précision@3, MRR@3 et nDCG@3 du moteur hybride sous annotation
manuelle & Annotation, 40 requêtes & \texttt{rag\_\allowbreak protocoles} \\
\hline
0,4396 / 0,4604 & Rappel@3 de l'hybride et du sémantique, que seule l'annotation rend
calculable & Annotation & idem \\
\hline
22 / 152 & Couples (requête, résolution) jugés utiles sur les couples distincts &
Annotation & \texttt{feuille\_\allowbreak jugement} \\
\hline
\end{longtable}

\textbf{Une exception, que nous signalons plutôt que de la masquer.} Les quatre valeurs du
tableau~\ref{tab:baseline} — 0,3746 et 0,4789 notamment — proviennent du carnet d'expérimentation
\texttt{05\_camembert\_finetuning}, où elles sont inscrites en dur dans une cellule de synthèse.
Elles ont été mesurées, mais aucun artefact du dépôt ne permet de les recalculer, à la différence
de toutes les autres lignes de ce tableau. Elles servent à montrer que quatre familles de modèles
s'arrêtent au même niveau, constat que la campagne tracée confirme par ailleurs avec 0,3753 pour le
meilleur d'entre eux.

\FloatBarrier
